%% Fasit til kortsett fra mattelab.no – laget automatisk av bygg.py.
%% Kompiler med lualatex.
\documentclass[11pt,a4paper]{article}
\usepackage[norsk]{babel}
\usepackage[a4paper,margin=1.2cm,bottom=1.7cm]{geometry}
\usepackage{fontspec}
\usepackage{unicode-math}
\setmainfont{TeX Gyre Pagella}
\setsansfont{TeX Gyre Heros}[Scale=MatchLowercase]
\setmathfont{TeX Gyre Pagella Math}
\usepackage{xcolor}
\definecolor{blekk}{HTML}{16213A}
\definecolor{dempet}{HTML}{566077}
\definecolor{m1}{HTML}{2F4AB0}
\usepackage{tikz}
\usepackage{enumitem}
\usepackage{qrcode}
\usepackage{fancyhdr}
\usepackage[hidelinks]{hyperref}
\pdfvariable suppressoptionalinfo \numexpr 1+2+4+8+16+32+64+128+256+512\relax
\setlength{\parindent}{0pt}
\setlength{\parskip}{0pt}
\pagestyle{fancy}
\fancyhf{}
\renewcommand{\headrulewidth}{0pt}
\fancyfoot[L]{\sffamily\footnotesize\color{dempet}mattelab.no · Potenser i rekkefølge · kortsett K4 · fasit}
\fancyfoot[R]{\sffamily\footnotesize\color{dempet}Side \thepage}
\begin{document}
\begin{minipage}[t]{0.8\linewidth}\vspace{0pt}
{\sffamily\small\color{dempet}Fasit for læreren · kortsett K4 · 1T}\par\vspace{0.25em}
{\fontsize{20}{24}\selectfont\bfseries Potenser i rekkefølge}\par\vspace{0.5em}
Legg kortene i rekkefølge fra minst til størst. Regn uten kalkulator.\par\vspace{0.4em}{\small\color{dempet}Kortene er nummerert i den rekkefølgen de står i utskriften.}
\end{minipage}\hfill
\begin{minipage}[t]{0.15\linewidth}\vspace{0pt}\raggedleft
\qrcode[height=1.9cm]{https://mattelab.no/varierende/potenser-i-rekkefolge.html}
\end{minipage}\par\vspace{0.8em}

\begin{enumerate}[leftmargin=1.6em, itemsep=0.45em]
\item \textsf{\small K4-05}\quad $(-2)^3$\quad {\color{dempet}$= -8$}
\item \textsf{\small K4-03}\quad $-2^2$\quad {\color{dempet}$= -4$}
\item \textsf{\small K4-04}\quad $10^{-1}$\quad {\color{dempet}$= 0{,}1$}
\item \textsf{\small K4-02}\quad $2^{-2}$\quad {\color{dempet}$= 0{,}25$}
\item \textsf{\small K4-08}\quad $3^{-1}$\quad {\color{dempet}$= \tfrac{1}{3} \approx 0{,}333$}
\item \textsf{\small K4-10}\quad $5^0$\quad {\color{dempet}$= 1$}
\item \textsf{\small K4-09}\quad $0{,}5^{-1}$\quad {\color{dempet}$= 2$}
\item \textsf{\small K4-06}\quad $1{,}5^2$\quad {\color{dempet}$= 2{,}25$}
\item \textsf{\small K4-01}\quad $\left(\dfrac{1}{2}\right)^{-3}$\quad {\color{dempet}$= 8$}
\item \textsf{\small K4-07}\quad $3^2$\quad {\color{dempet}$= 9$}
\end{enumerate}
\subsection*{Til læreren}
Mange tror at en negativ eksponent gir et negativt tall. Kortene $2^{-2}$, $3^{-1}$ og $10^{-1}$ er positive, men mindre enn 1. Kortene $0{,}5^{-1}$ og $(1/2)^{-3}$ er større enn 1. Spør: Når blir en potens negativ?
\end{document}
